\documentclass[12pt,a4paper]{article}

\usepackage[a4paper,margin=1in]{geometry}
\usepackage{amsmath,amssymb,bm}
\usepackage{array}
\usepackage{mathtools}

\setlength{\parindent}{0pt}
\setlength{\parskip}{6pt}

\title{\textbf{Derivation of State-Space Matrix for MPC Control}}
\author{}
\date{}

\begin{document}

\maketitle

The following derivation follows the sequence and notation of the
handwritten derivation. The factor of $2$ is retained consistently in
the front and rear lateral tire forces, and the required signs in the
lateral-velocity and yaw-rate equations are corrected.

%====================================================================
\section{Vehicle Kinematic Equations}

The vehicle states are the global longitudinal position $X$, lateral
position $Y$, yaw angle $\psi$, longitudinal velocity $v_x$, lateral
velocity $v_y$, and yaw rate $r$.

The global longitudinal-position equation is

\begin{equation}
\boxed{
\dot X=v_x\cos\psi-v_y\sin\psi
}.
\end{equation}

The global lateral-position equation is

\begin{equation}
\boxed{
\dot Y=v_x\sin\psi+v_y\cos\psi
}.
\end{equation}

The yaw-rate relation is

\begin{equation}
\boxed{
\dot\psi=r
}.
\end{equation}

The longitudinal vehicle dynamics used in the handwritten model are

\begin{equation}
\boxed{
\dot v_x
=
a_{cmd}
-\mu_{rr}g\cos\theta
-g\sin\theta
+v_yr
}.
\end{equation}

Here, $\mu_{rr}g\cos\theta$ represents the rolling/frictional resistance,
$g\sin\theta$ represents the road-grade contribution, and $v_yr$ is the
body-frame coupling caused by yaw motion.

%====================================================================
\section{Front and Rear Normal Loads}

The front normal force is

\begin{equation}
\boxed{
F_{zf}
=
\frac{
mg\cos\theta\,a_2
-
mh\left(\dot v_x+g\sin\theta\right)
}
{a_1+a_2}
}.
\end{equation}

Similarly, the rear normal force is

\begin{equation}
\boxed{
F_{zr}
=
\frac{
mg\cos\theta\,a_1
+
mh\left(\dot v_x+g\sin\theta\right)
}
{a_1+a_2}
}.
\end{equation}

For the nominal condition,

\begin{equation}
F_{zf,\mathrm{nom}}
=
\frac{mga_2}{a_1+a_2},
\end{equation}

and

\begin{equation}
F_{zr,\mathrm{nom}}
=
\frac{mga_1}{a_1+a_2}.
\end{equation}

The front cornering stiffness is

\begin{equation}
C_{\alpha f}
=
C_{\alpha f,\mathrm{nom}}
\left(
\frac{F_{zf}}
{F_{zf,\mathrm{nom}}}
\right).
\end{equation}

Similarly,

\begin{equation}
C_{\alpha r}
=
C_{\alpha r,\mathrm{nom}}
\left(
\frac{F_{zr}}
{F_{zr,\mathrm{nom}}}
\right).
\end{equation}

The nominal values appearing in the handwritten derivation are

\begin{equation}
C_{\alpha f,\mathrm{nom}}
=
57{,}000\;\mathrm{N/rad},
\end{equation}

and

\begin{equation}
C_{\alpha r,\mathrm{nom}}
=
52{,}000\;\mathrm{N/rad}.
\end{equation}

%====================================================================
\section{Front and Rear Slip Angles}

The front slip angle is

\begin{equation}
\boxed{
\alpha_f
=
\delta
-
\frac{v_y+a_1r}{v_x}
}.
\end{equation}

The rear slip angle is

\begin{equation}
\boxed{
\alpha_r
=
-\frac{v_y-a_2r}{v_x}
=
\frac{-v_y+a_2r}{v_x}
}.
\end{equation}

%====================================================================
\section{Lateral Tire Forces}

Since the cornering stiffnesses are defined for individual tires, two
tires contribute at each axle.

Therefore,

\begin{equation}
\boxed{
F_{yf}=2C_{\alpha f}\alpha_f
},
\end{equation}

and

\begin{equation}
\boxed{
F_{yr}=2C_{\alpha r}\alpha_r
}.
\end{equation}

%====================================================================
\section{Newton--Euler Equations in the Vehicle-Fixed Frame}

Applying Newton's second law in the lateral direction,

\begin{equation}
\boxed{
m\left(\dot v_y+v_xr\right)
=
F_{yf}+F_{yr}
}.
\end{equation}

Applying the yaw-moment equation about the center of gravity,

\begin{equation}
\boxed{
I_z\dot r
=
a_1F_{yf}-a_2F_{yr}
}.
\end{equation}

%====================================================================
\subsection{Detailed Derivation of the Lateral-Velocity Equation}

Starting from

\begin{equation}
m\left(\dot v_y+v_xr\right)
=
F_{yf}+F_{yr},
\end{equation}

substitute

\begin{equation}
F_{yf}=2C_{\alpha f}\alpha_f,
\qquad
F_{yr}=2C_{\alpha r}\alpha_r.
\end{equation}

Thus,

\begin{equation}
m\left(\dot v_y+v_xr\right)
=
2C_{\alpha f}\alpha_f
+
2C_{\alpha r}\alpha_r.
\end{equation}

Expanding the left-hand side,

\begin{equation}
m\dot v_y
+
mv_xr
=
2C_{\alpha f}\alpha_f
+
2C_{\alpha r}\alpha_r.
\end{equation}

Moving $mv_xr$ to the right-hand side,

\begin{equation}
m\dot v_y
=
2C_{\alpha f}\alpha_f
+
2C_{\alpha r}\alpha_r
-
mv_xr.
\end{equation}

Dividing by $m$,

\begin{equation}
\dot v_y
=
\frac{2C_{\alpha f}}{m}\alpha_f
+
\frac{2C_{\alpha r}}{m}\alpha_r
-
v_xr.
\end{equation}

Now,

\begin{equation}
\alpha_f
=
\delta
-
\frac{v_y+a_1r}{v_x},
\end{equation}

and

\begin{equation}
\alpha_r
=
\frac{-v_y+a_2r}{v_x}.
\end{equation}

Substituting these equations,

\begin{align}
\dot v_y
={}&
\frac{2C_{\alpha f}}{m}
\left(
\delta-\frac{v_y+a_1r}{v_x}
\right)
\nonumber\\
&+
\frac{2C_{\alpha r}}{m}
\left(
\frac{-v_y+a_2r}{v_x}
\right)
-v_xr.
\end{align}

Expanding the front term,

\begin{align}
\frac{2C_{\alpha f}}{m}
\left(
\delta-\frac{v_y+a_1r}{v_x}
\right)
={}&
\frac{2C_{\alpha f}}{m}\delta
-
\frac{2C_{\alpha f}}{mv_x}v_y
\nonumber\\
&-
\frac{2a_1C_{\alpha f}}{mv_x}r.
\end{align}

Expanding the rear term,

\begin{equation}
\frac{2C_{\alpha r}}{m}
\left(
\frac{-v_y+a_2r}{v_x}
\right)
=
-\frac{2C_{\alpha r}}{mv_x}v_y
+
\frac{2a_2C_{\alpha r}}{mv_x}r.
\end{equation}

Therefore,

\begin{align}
\dot v_y
={}&
\frac{2C_{\alpha f}}{m}\delta
-\frac{2C_{\alpha f}}{mv_x}v_y
-\frac{2a_1C_{\alpha f}}{mv_x}r
\nonumber\\
&-\frac{2C_{\alpha r}}{mv_x}v_y
+\frac{2a_2C_{\alpha r}}{mv_x}r
-v_xr.
\end{align}

Collecting the $v_y$ terms,

\begin{align}
&
-\frac{2C_{\alpha f}}{mv_x}v_y
-
\frac{2C_{\alpha r}}{mv_x}v_y
\nonumber\\
&=
-v_y
\left(
\frac{2C_{\alpha f}+2C_{\alpha r}}{mv_x}
\right).
\end{align}

Collecting the $r$ terms,

\begin{align}
&
-\frac{2a_1C_{\alpha f}}{mv_x}r
+
\frac{2a_2C_{\alpha r}}{mv_x}r
-v_xr
\nonumber\\
&=
r
\left(
\frac{-2a_1C_{\alpha f}+2a_2C_{\alpha r}}
{mv_x}
-v_x
\right).
\end{align}

Therefore,

\begin{equation}
\boxed{
\dot v_y
=
-v_y
\left(
\frac{2C_{\alpha f}+2C_{\alpha r}}
{mv_x}
\right)
+
r
\left(
\frac{-2a_1C_{\alpha f}+2a_2C_{\alpha r}}
{mv_x}
-v_x
\right)
+
\frac{2C_{\alpha f}}{m}\delta
}.
\end{equation}

The front-axle yaw-rate contribution has a negative sign because the front
slip-angle equation contains

\begin{equation}
-\frac{a_1r}{v_x}.
\end{equation}

%====================================================================
\subsection{Detailed Derivation of the Yaw-Rate Equation}

Starting from

\begin{equation}
I_z\dot r
=
a_1F_{yf}
-
a_2F_{yr},
\end{equation}

substitute

\begin{equation}
F_{yf}=2C_{\alpha f}\alpha_f,
\end{equation}

and

\begin{equation}
F_{yr}=2C_{\alpha r}\alpha_r.
\end{equation}

Hence,

\begin{equation}
I_z\dot r
=
a_1\left(2C_{\alpha f}\alpha_f\right)
-
a_2\left(2C_{\alpha r}\alpha_r\right).
\end{equation}

Therefore,

\begin{equation}
I_z\dot r
=
2a_1C_{\alpha f}\alpha_f
-
2a_2C_{\alpha r}\alpha_r.
\end{equation}

Substituting the front and rear slip angles,

\begin{align}
I_z\dot r
={}&
2a_1C_{\alpha f}
\left(
\delta-\frac{v_y+a_1r}{v_x}
\right)
\nonumber\\
&-
2a_2C_{\alpha r}
\left(
\frac{-v_y+a_2r}{v_x}
\right).
\end{align}

Expanding the front contribution,

\begin{align}
2a_1C_{\alpha f}
\left(
\delta-\frac{v_y+a_1r}{v_x}
\right)
={}&
2a_1C_{\alpha f}\delta
-
\frac{2a_1C_{\alpha f}}{v_x}v_y
\nonumber\\
&-
\frac{2a_1^2C_{\alpha f}}{v_x}r.
\end{align}

Expanding the rear contribution,

\begin{equation}
-2a_2C_{\alpha r}
\left(
\frac{-v_y+a_2r}{v_x}
\right)
=
+
\frac{2a_2C_{\alpha r}}{v_x}v_y
-
\frac{2a_2^2C_{\alpha r}}{v_x}r.
\end{equation}

Hence,

\begin{align}
I_z\dot r
={}&
2a_1C_{\alpha f}\delta
-
\frac{2a_1C_{\alpha f}}{v_x}v_y
-
\frac{2a_1^2C_{\alpha f}}{v_x}r
\nonumber\\
&+
\frac{2a_2C_{\alpha r}}{v_x}v_y
-
\frac{2a_2^2C_{\alpha r}}{v_x}r.
\end{align}

Collecting the $v_y$ terms,

\begin{equation}
-\frac{2a_1C_{\alpha f}}{v_x}v_y
+
\frac{2a_2C_{\alpha r}}{v_x}v_y
=
\frac{-2a_1C_{\alpha f}+2a_2C_{\alpha r}}
{v_x}v_y.
\end{equation}

Collecting the $r$ terms,

\begin{align}
&
-\frac{2a_1^2C_{\alpha f}}{v_x}r
-
\frac{2a_2^2C_{\alpha r}}{v_x}r
\nonumber\\
&=
-\frac{
2a_1^2C_{\alpha f}
+
2a_2^2C_{\alpha r}
}{v_x}r.
\end{align}

Therefore,

\begin{align}
I_z\dot r
={}&
2a_1C_{\alpha f}\delta
+
\frac{-2a_1C_{\alpha f}+2a_2C_{\alpha r}}
{v_x}v_y
\nonumber\\
&-
\frac{
2a_1^2C_{\alpha f}
+
2a_2^2C_{\alpha r}
}{v_x}r.
\end{align}

Dividing throughout by $I_z$,

\begin{equation}
\boxed{
\dot r
=
\frac{2a_1C_{\alpha f}}{I_z}\delta
+
\frac{-2a_1C_{\alpha f}+2a_2C_{\alpha r}}
{I_zv_x}v_y
-
\frac{
2a_1^2C_{\alpha f}
+
2a_2^2C_{\alpha r}
}
{I_zv_x}r
}.
\end{equation}

%====================================================================
\section{Continuous State-Space Representation}

Using the structure shown in the handwritten derivation,

\begin{equation}
\boxed{
\dot{\bm{x}}
=
A\bm{x}
+
B\bm{u}
+
B_{ud}\bm{d}_{um}
+
B_{md}\bm{d}_m
}.
\end{equation}

The output equation is

\begin{equation}
\boxed{
\bm{y}
=
C\bm{x}
+
D\bm{u}
}.
\end{equation}

The state vector is

\begin{equation}
\boxed{
\bm{x}
=
\begin{bmatrix}
X\\
Y\\
\psi\\
v_x\\
v_y\\
r
\end{bmatrix}
}.
\end{equation}

The control input is

\begin{equation}
\boxed{
\bm{u}
=
\begin{bmatrix}
\delta\\
a_{cmd}
\end{bmatrix}
}.
\end{equation}

For linearization around an operating point $(\bm{x}_0,\bm{u}_0)$,

\begin{equation}
\boxed{
A
=
\left.
\frac{\partial f}{\partial\bm{x}}
\right|_{\bm{x}_0,\bm{u}_0}
},
\end{equation}

and

\begin{equation}
\boxed{
B
=
\left.
\frac{\partial f}{\partial\bm{u}}
\right|_{\bm{x}_0,\bm{u}_0}
}.
\end{equation}

%====================================================================
\section{Derivation of the $A$ Matrix}

From

\begin{equation}
\dot X
=
v_x\cos\psi
-
v_y\sin\psi,
\end{equation}

the nonzero partial derivatives are

\begin{equation}
\frac{\partial\dot X}{\partial\psi}
=
-v_x\sin\psi
-v_y\cos\psi,
\end{equation}

\begin{equation}
\frac{\partial\dot X}{\partial v_x}
=
\cos\psi,
\end{equation}

and

\begin{equation}
\frac{\partial\dot X}{\partial v_y}
=
-\sin\psi.
\end{equation}

Hence,

\begin{equation}
A_{13}
=
-v_{x0}\sin\psi_0
-v_{y0}\cos\psi_0,
\end{equation}

\begin{equation}
A_{14}
=
\cos\psi_0,
\end{equation}

and

\begin{equation}
A_{15}
=
-\sin\psi_0.
\end{equation}

From

\begin{equation}
\dot Y
=
v_x\sin\psi
+
v_y\cos\psi,
\end{equation}

we obtain

\begin{equation}
\frac{\partial\dot Y}{\partial\psi}
=
v_x\cos\psi
-v_y\sin\psi,
\end{equation}

\begin{equation}
\frac{\partial\dot Y}{\partial v_x}
=
\sin\psi,
\end{equation}

and

\begin{equation}
\frac{\partial\dot Y}{\partial v_y}
=
\cos\psi.
\end{equation}

Hence,

\begin{equation}
A_{23}
=
v_{x0}\cos\psi_0
-v_{y0}\sin\psi_0,
\end{equation}

\begin{equation}
A_{24}
=
\sin\psi_0,
\end{equation}

and

\begin{equation}
A_{25}
=
\cos\psi_0.
\end{equation}

Since

\begin{equation}
\dot\psi=r,
\end{equation}

we obtain

\begin{equation}
\boxed{
A_{36}=1
}.
\end{equation}

For

\begin{equation}
\dot v_x
=
a_{cmd}
-\mu_{rr}g\cos\theta
-g\sin\theta
+v_yr,
\end{equation}

the relevant state derivatives are

\begin{equation}
\frac{\partial\dot v_x}{\partial v_y}
=
r,
\end{equation}

and

\begin{equation}
\frac{\partial\dot v_x}{\partial r}
=
v_y.
\end{equation}

Thus,

\begin{equation}
A_{45}=r_0,
\end{equation}

and

\begin{equation}
A_{46}=v_{y0}.
\end{equation}

For the lateral equation,

\begin{align}
\dot v_y
={}&
-\frac{2C_{\alpha f}+2C_{\alpha r}}
{mv_x}v_y
\nonumber\\
&+
\left(
\frac{-2a_1C_{\alpha f}+2a_2C_{\alpha r}}
{mv_x}
-v_x
\right)r
+
\frac{2C_{\alpha f}}{m}\delta.
\end{align}

Differentiating with respect to $v_x$,

\begin{align}
\frac{\partial\dot v_y}{\partial v_x}
={}&
\frac{
(2C_{\alpha f}+2C_{\alpha r})v_y
}
{mv_x^2}
\nonumber\\
&-
\frac{
(-2a_1C_{\alpha f}+2a_2C_{\alpha r})r
}
{mv_x^2}
-r.
\end{align}

Therefore,

\begin{equation}
\boxed{
A_{54}
=
\frac{
(2C_{\alpha f}+2C_{\alpha r})v_{y0}
-
(-2a_1C_{\alpha f}+2a_2C_{\alpha r})r_0
}
{mv_{x0}^2}
-r_0
}.
\end{equation}

Differentiating with respect to $v_y$,

\begin{equation}
\boxed{
A_{55}
=
-\frac{
2C_{\alpha f}+2C_{\alpha r}
}
{mv_{x0}}
}.
\end{equation}

Differentiating with respect to $r$,

\begin{equation}
\boxed{
A_{56}
=
\frac{
-2a_1C_{\alpha f}
+
2a_2C_{\alpha r}
}
{mv_{x0}}
-v_{x0}
}.
\end{equation}

For the yaw-rate equation,

\begin{align}
\dot r
={}&
\frac{
-2a_1C_{\alpha f}
+
2a_2C_{\alpha r}
}
{I_zv_x}v_y
\nonumber\\
&-
\frac{
2a_1^2C_{\alpha f}
+
2a_2^2C_{\alpha r}
}
{I_zv_x}r
+
\frac{2a_1C_{\alpha f}}{I_z}\delta.
\end{align}

Differentiating with respect to $v_x$,

\begin{align}
\frac{\partial\dot r}{\partial v_x}
={}&
-\frac{
(-2a_1C_{\alpha f}+2a_2C_{\alpha r})v_y
}
{I_zv_x^2}
\nonumber\\
&+
\frac{
(2a_1^2C_{\alpha f}+2a_2^2C_{\alpha r})r
}
{I_zv_x^2}.
\end{align}

Thus,

\begin{equation}
\boxed{
A_{64}
=
\frac{
-(-2a_1C_{\alpha f}+2a_2C_{\alpha r})v_{y0}
+
(2a_1^2C_{\alpha f}+2a_2^2C_{\alpha r})r_0
}
{I_zv_{x0}^2}
}.
\end{equation}

Differentiating with respect to $v_y$,

\begin{equation}
\boxed{
A_{65}
=
\frac{
-2a_1C_{\alpha f}
+
2a_2C_{\alpha r}
}
{I_zv_{x0}}
}.
\end{equation}

Differentiating with respect to $r$,

\begin{equation}
\boxed{
A_{66}
=
-\frac{
2a_1^2C_{\alpha f}
+
2a_2^2C_{\alpha r}
}
{I_zv_{x0}}
}.
\end{equation}

Therefore,

\begin{equation}
\boxed{
A=
\begin{bmatrix}
0 &
0 &
-v_{x0}\sin\psi_0-v_{y0}\cos\psi_0 &
\cos\psi_0 &
-\sin\psi_0 &
0
\\[1mm]
0 &
0 &
v_{x0}\cos\psi_0-v_{y0}\sin\psi_0 &
\sin\psi_0 &
\cos\psi_0 &
0
\\[1mm]
0&0&0&0&0&1
\\[1mm]
0&0&0&0&r_0&v_{y0}
\\[1mm]
0&0&0&A_{54}&
-\dfrac{2C_{\alpha f}+2C_{\alpha r}}{mv_{x0}}&
\dfrac{-2a_1C_{\alpha f}+2a_2C_{\alpha r}}{mv_{x0}}
-v_{x0}
\\[4mm]
0&0&0&A_{64}&
\dfrac{-2a_1C_{\alpha f}+2a_2C_{\alpha r}}{I_zv_{x0}}&
-\dfrac{2a_1^2C_{\alpha f}+2a_2^2C_{\alpha r}}
{I_zv_{x0}}
\end{bmatrix}
}.
\end{equation}

%====================================================================
\section{Derivation of the $B$ Matrix}

The input vector is

\begin{equation}
\bm{u}
=
\begin{bmatrix}
\delta\\
a_{cmd}
\end{bmatrix}.
\end{equation}

The steering angle $\delta$ enters the lateral-velocity equation through

\begin{equation}
\frac{2C_{\alpha f}}{m}\delta.
\end{equation}

Therefore,

\begin{equation}
\frac{\partial\dot v_y}{\partial\delta}
=
\frac{2C_{\alpha f}}{m}.
\end{equation}

Similarly, the steering angle enters the yaw-rate equation through

\begin{equation}
\frac{2a_1C_{\alpha f}}{I_z}\delta.
\end{equation}

Therefore,

\begin{equation}
\frac{\partial\dot r}{\partial\delta}
=
\frac{2a_1C_{\alpha f}}{I_z}.
\end{equation}

The longitudinal acceleration command appears in

\begin{equation}
\dot v_x
=
a_{cmd}
-\mu_{rr}g\cos\theta
-g\sin\theta
+v_yr.
\end{equation}

Therefore,

\begin{equation}
\frac{\partial\dot v_x}{\partial a_{cmd}}
=
1.
\end{equation}

Hence,

\begin{equation}
\boxed{
B
=
\begin{bmatrix}
0&0\\
0&0\\
0&0\\
0&1\\
\dfrac{2C_{\alpha f}}{m}&0\\
\dfrac{2a_1C_{\alpha f}}{I_z}&0
\end{bmatrix}
}.
\end{equation}

%====================================================================
\section{Unmodelled-Force Disturbance Matrix}

The unmodelled-force disturbance vector is

\begin{equation}
\boxed{
\bm{d}_{um}
=
\begin{bmatrix}
\Delta F_{x,\mathrm{frict}}\\
\Delta F_{yf}\\
\Delta F_{yr}
\end{bmatrix}
}.
\end{equation}

The longitudinal-force disturbance contributes

\begin{equation}
\Delta\dot v_x
=
\frac{\Delta F_{x,\mathrm{frict}}}{m}.
\end{equation}

The lateral-force disturbances contribute

\begin{equation}
\Delta\dot v_y
=
\frac{
\Delta F_{yf}
+
\Delta F_{yr}
}{m}.
\end{equation}

The yaw-rate disturbance contribution is

\begin{equation}
\Delta\dot r
=
\frac{
a_1\Delta F_{yf}
-
a_2\Delta F_{yr}
}{I_z}.
\end{equation}

Therefore,

\begin{equation}
\boxed{
B_{ud}
=
\begin{bmatrix}
0&0&0\\
0&0&0\\
0&0&0\\
\dfrac{1}{m}&0&0\\
0&\dfrac{1}{m}&\dfrac{1}{m}\\
0&\dfrac{a_1}{I_z}&-\dfrac{a_2}{I_z}
\end{bmatrix}
}.
\end{equation}

The modeled disturbance vector in the handwritten notes is

\begin{equation}
\boxed{
\bm{d}_m
=
\begin{bmatrix}
\delta_{\mathrm{user}}\\
\theta_{\mathrm{grade}}
\end{bmatrix}
}.
\end{equation}

Accordingly, $B_{md}$ represents the corresponding modeled-disturbance
input matrix.

%====================================================================
\section{Output Equation}

The output matrix selects $X$, $Y$, and $\psi$:

\begin{equation}
\boxed{
C
=
\begin{bmatrix}
1&0&0&0&0&0\\
0&1&0&0&0&0\\
0&0&1&0&0&0
\end{bmatrix}
}.
\end{equation}

For zero direct feedthrough,

\begin{equation}
\boxed{
D=0
}.
\end{equation}

Hence,

\begin{equation}
\boxed{
\bm{y}
=
C\bm{x}
+
D\bm{u}
=
C\bm{x}
}.
\end{equation}

%====================================================================
\section{Final Corrected Continuous Model}

The complete corrected continuous vehicle model is

\begin{equation}
\boxed{
\begin{aligned}
\dot X
&=
v_x\cos\psi-v_y\sin\psi,
\\[2mm]
\dot Y
&=
v_x\sin\psi+v_y\cos\psi,
\\[2mm]
\dot\psi
&=
r,
\\[2mm]
\dot v_x
&=
a_{cmd}
-\mu_{rr}g\cos\theta
-g\sin\theta
+v_yr,
\\[2mm]
\dot v_y
&=
-v_y
\left(
\frac{2C_{\alpha f}+2C_{\alpha r}}
{mv_x}
\right)
+
r
\left(
\frac{-2a_1C_{\alpha f}+2a_2C_{\alpha r}}
{mv_x}
-v_x
\right)
+
\frac{2C_{\alpha f}}{m}\delta,
\\[2mm]
\dot r
&=
\frac{-2a_1C_{\alpha f}+2a_2C_{\alpha r}}
{I_zv_x}v_y
-
\frac{
2a_1^2C_{\alpha f}
+
2a_2^2C_{\alpha r}
}
{I_zv_x}r
+
\frac{2a_1C_{\alpha f}}{I_z}\delta.
\end{aligned}
}
\end{equation}

\end{document}